% Generated draft: requires mathematical and source review.
\documentclass[11pt]{article}
\usepackage[margin=25mm]{geometry}
\usepackage{fontspec}
\setmainfont{Latin Modern Roman}
\usepackage{amsmath,amssymb,mathtools,unicode-math}
\setmathfont{Latin Modern Math}
\usepackage{enumitem,needspace,xurl,hyperref}
\setlength{\parindent}{0pt}
\setlength{\parskip}{5pt}
\setlength{\emergencystretch}{3em}
\newcommand{\N}{\mathbb{N}}
\newcommand{\Z}{\mathbb{Z}}
\newcommand{\Q}{\mathbb{Q}}
\newcommand{\R}{\mathbb{R}}
\newcommand{\C}{\mathbb{C}}
\newcommand{\F}{\mathbb{F}}
\newcommand{\A}{\mathbb{A}}
\newcommand{\PP}{\mathbb{P}}
\newcommand{\E}{\mathbb{E}}
\newcommand{\eps}{\varepsilon}
\newcommand{\dd}{\,\mathrm{d}}
\begin{document}
\section*{UnsolvedMath ID: 2789}
\textbf{Source problem code:} catalogue entry
\subsection{Definitions and mathematical statement}
An infinite-type surface is a surface not homeomorphic to a compact surface with finitely many points removed. The source asks for a fixed finite list of practically computable group invariants that recovers the homeomorphism type of such a surface from either its mapping class group or its pure mapping class group.

In the usual orientable convention, \(\operatorname{Map}(S)\) consists of isotopy classes of orientation-preserving homeomorphisms, with boundary conditions fixed as part of the convention. The pure group is the subgroup fixing every end individually. The source does not explicitly specify orientability, boundary treatment, full versus orientation-preserving groups, or all infinite-type hypotheses, so these contextual definitions do not silently restrict the intended surface class.

For a specified class \(\mathcal S\) and one chosen group assignment \(H(S)=\operatorname{Map}(S)\) or \(H(S)=\operatorname{PMap}(S)\), the desired separation property is a finite collection \(I_1,\ldots,I_r\) such that
\[
\forall S,S'\in\mathcal S:\qquad
\left[\forall j\in\{1,\ldots,r\},\ I_j(H(S))=I_j(H(S'))\right]
\Longrightarrow S\cong S'.
\]
Each \(I_j\) must depend only on the intended group-isomorphism type, and the list is independent of \(S\). The essential additional demand is practical computation of its values. The source gives no input presentation, oracle model, permitted output data or formal complexity criterion for these generally uncountable groups. The bare separation implication alone could be met tautologically by full group-isomorphism data and therefore does not capture the practical requirement. No such tautological invariant is proposed as a solution.

The supplied background describes infinite-type classification through end space \(E(S)\), the closed subset \(E^g(S)\) of ends accumulated by genus, genus and compact boundary count. It calls \(E^g\) a closed subgroup, but supplies no group law on the end space; the intended topological object is a closed subset. The full orientation and boundary hypotheses needed for the cited classification are not reproduced. It stresses the practical difficulty of comparing closed subsets of a Cantor set and records Bavard--Dowdall--Rafi's algebraic rigidity statement that isomorphic big mapping class groups determine homeomorphic surfaces, while noting that this does not make isomorphism testing practical.

For finite-type context, the source cites Birman--Lubotzky--McCarthy's algebraic-rank work and Harer's virtual cohomological dimension, two exceptional small-surface pairs, and a limitation when \(g=0\) or \(b=0\), referring to Rafi--Schleimer, Remark A.2. It reports one result in the desired direction from Aougab--Patel--Vlamis, Theorem 8.1, confined to orientable infinite-genus surfaces with finitely many ends and no planar ends. Neither the exceptional pairs nor the exact invariant list in that restricted theorem is supplied, so neither is guessed.

The registry and supplied triage dated 2026-08-17 label the practical finite-list reconstruction open. The triage reports classification data as partial context, not a finite practically computable invariant list or an impossibility theorem.
\subsection{Short English statement}
Find finitely many usable invariants of an infinite-type surface's mapping class group, or of its pure subgroup, that distinguish its topology. Algebraic rigidity alone does not provide practical computation. The source leaves the surface conventions and computation model unspecified, and its known directional result has a restricted finite-ended infinite-genus scope.
\subsection{Sources}
\begin{itemize}
\item Local UnsolvedMath registry, permanent ID 2789, source code catalogue entry, ``Kirby Problem 2.41,'' including supplied literature triage dated 2026-08-17.
\item \emph{K3: A New Problem List in Low-Dimensional Topology}, Problem 2.41 (2026), supplied source: \url{https://aimath.org/pastworkshops/kirbylistrep.pdf}.
\item Joan S. Birman, Alex Lubotzky, John McCarthy, et al., as listed in the supplied bibliography, ``Abelian and solvable subgroups of the mapping class groups,'' \emph{Duke Mathematical Journal} 50(4) (1983), 1107--1120, source label [BLM+83].
\item John L. Harer, ``The virtual cohomological dimension of the mapping class group of an orientable surface,'' \emph{Inventiones Mathematicae} 84(1) (1986), 157--176, DOI 10.1007/BF01388737, source label [Har86].
\item Kasra Rafi and Saul Schleimer, ``Curve complexes are rigid,'' \emph{Duke Mathematical Journal} 158(2) (2011), 225--246, DOI 10.1215/00127094-1334004, source label [RS11], especially the supplied pointer to Remark A.2.
\item Juliette Bavard, Spencer Dowdall, and Kasra Rafi, ``Isomorphisms between big mapping class groups,'' \emph{International Mathematics Research Notices} 2020(10) (2020), 3084--3099, source label [BDR20].
\item Tarik Aougab, Priyam Patel, and Nicholas G. Vlamis, ``Isometry groups of infinite genus hyperbolic surfaces,'' \emph{Mathematische Annalen} 381 (2021), 459--498, source label [APV21], Theorem 8.1.
\end{itemize}
\end{document}
